\documentclass[10pt]{article}
\usepackage[a4paper,margin=0.86in]{geometry}
\usepackage[T1]{fontenc}
\usepackage{lmodern}
\usepackage{amsmath,amssymb,amsthm}
\usepackage{graphicx}
\usepackage{microtype}
\usepackage[colorlinks=true,linkcolor=blue,urlcolor=blue,citecolor=blue]{hyperref}
\newtheorem{theorem}{Theorem}
\newtheorem{lemma}{Lemma}
\newcommand{\M}{\mathcal M}
\newcommand{\N}{\mathbb N}
\newcommand{\Tr}{\operatorname{Tr}}
\title{The Grothendieck Criterion Extends to\\
Bounded-Degree Type-I Marcinkiewicz Spaces}
\author{Candidate strong partial-result packet\\
\small arXiv:1912.01162, Introduction}
\date{12 August 2026}

\begin{document}
\maketitle

\paragraph{Status.}
Candidate strong partial answer, subject to expert review.  The scalar
criterion extends exactly to every fixed finite matrix amplification and,
more generally, the positive result extends to all semifinite type-I von
Neumann algebras of uniformly bounded degree.  The source's unrestricted
noncommutative question remains open here.

\section{Source question}
De Pagter--Sukochev characterize the increasing concave functions $\psi$ for
which the scalar Marcinkiewicz space $M_\psi(0,\gamma)$ is Grothendieck.  In
their Introduction, page 2, they ask:
\begin{quote}
``It seems to be an interesting open question whether the results presented
in this article also hold for noncommutative Marcinkiewicz spaces.''
\end{quote}
Their conditions are
\begin{align*}
\text{(A)}\quad &\gamma=\infty,
 &&\liminf_{t\downarrow0}\frac{\psi(2t)}{\psi(t)}>1,
 &&\liminf_{t\to\infty}\frac{\psi(2t)}{\psi(t)}>1;\\
\text{(B)}\quad &\gamma<\infty,
 &&\liminf_{t\downarrow0}\frac{\psi(2t)}{\psi(t)}>1.
\end{align*}
They prove that (A) or (B), as appropriate, is equivalent to
$M_\psi(0,\gamma)$ being Grothendieck, and gives sufficiency over arbitrary
$\sigma$-finite scalar measure spaces \cite[Theorems 5.6 and 6.3]{DPS}.

\begin{figure}[ht]
\centering
\includegraphics[width=.96\textwidth]{figures/open_question_crop.png}
\caption{The open question in the source PDF, page 2.}
\end{figure}

\section{Definitions}
Let $(\M,\tau)$ be a semifinite von Neumann algebra with normal faithful
semifinite trace, and let $\mu(x)$ denote the generalized singular-value
function of a $\tau$-measurable operator $x$.  The noncommutative
Marcinkiewicz space is
\[
 M_\psi(\M,\tau)=\{x:\mu(x)\in M_\psi(0,\tau(1))\},\qquad
 \|x\|_{M_\psi(\M)}=\|\mu(x)\|_{M_\psi}.
\]
We call $(\M,\tau)$ a \emph{bounded-degree type-I model} if, for some
$d<\infty$,
\begin{equation}\label{eq:decomp}
 (\M,\tau)\cong
 \bigoplus_{r=1}^{d}
 \bigl(L_\infty(\Omega_r,\nu_r)\,\overline\otimes\,M_r,
       \int_{\Omega_r}\Tr_r\,d\nu_r\bigr).
\end{equation}
This is the standard central form of a semifinite subhomogeneous von Neumann
algebra.  Zero summands are allowed.

\section{Proof Intuition}
A fixed matrix algebra introduces only finitely many noncommutative
coordinates.  Compressing an operator-valued function by two matrix units
extracts one scalar entry and cannot increase any symmetric operator norm.
Conversely, summing the finitely many entry corners reconstructs the
operator.  Hence the noncommutative space is Banach-isomorphic to a finite
product of scalar Marcinkiewicz spaces.  Grothendieckness is stable under
finite products, and a scalar corner is complemented, giving the converse in
the homogeneous case.

\section{The matrix-corner lemma}
\begin{lemma}\label{lem:corners}
Let $E$ be any fully symmetric Banach function space and
$\M=L_\infty(\Omega)\overline\otimes M_r$ with the tensor trace.  For
$x=(x_{ij})\in E(\M)$,
\begin{equation}\label{eq:corners}
 \max_{i,j}\|x_{ij}e_{ij}\|_{E(\M)}
 \le \|x\|_{E(\M)}
 \le \sum_{i,j}\|x_{ij}e_{ij}\|_{E(\M)}
 \le r^2\max_{i,j}\|x_{ij}e_{ij}\|_{E(\M)}.
\end{equation}
Consequently, the entry map is a Banach isomorphism from $E(\M)$ onto the
finite product of its $r^2$ scalar corner spaces.
\end{lemma}

\begin{proof}
Let $e_{ij}$ be the constant matrix units.  The ideal property of a symmetric
operator space gives
\[
 \|e_{ii}xe_{jj}\|_{E(\M)}\le
 \|e_{ii}\|_\infty\,\|x\|_{E(\M)}\,\|e_{jj}\|_\infty
 \le \|x\|_{E(\M)}.
\]
But $e_{ii}xe_{jj}=x_{ij}e_{ij}$, proving the first inequality.  Since
$x=\sum_{i,j}x_{ij}e_{ij}$, the triangle inequality proves the second, and
the third is immediate.  Every family of measurable scalar entries defines
a measurable $M_r$-valued function; (\ref{eq:corners}) shows that it lies in
$E(\M)$ exactly when all its finitely many corners lie there.  Thus the entry
map and its inverse are bounded.
\end{proof}

For a corner over a measure space of mass $\gamma_r$, the singular-value
function of $fe_{ij}$ is the decreasing rearrangement of $|f|$, followed by
zeros.  Since $\psi$ is increasing, the Marcinkiewicz supremum after
$\gamma_r$ cannot exceed its value at $\gamma_r$.  Thus the corner space is
precisely the scalar $M_\psi(\Omega_r,\nu_r)$ with its usual norm.

\section{Partial solution}
\begin{theorem}\label{thm:main}
Let $\psi$ satisfy condition (A) when $\tau(1)=\infty$ and condition (B)
when $\tau(1)<\infty$.  For every bounded-degree type-I model
$(\M,\tau)$ in (\ref{eq:decomp}), the noncommutative Marcinkiewicz space
$M_\psi(\M,\tau)$ is a Grothendieck Banach space.

For the homogeneous atomless model
\[
 \M=L_\infty(0,\gamma)\overline\otimes M_r
 \quad (1\le r<\infty),
\]
the converse also holds: $M_\psi(\M,\tau)$ is Grothendieck if and only if
the applicable condition (A) or (B) holds.  In particular, every finite
matrix amplification of weak $L_p$, $1<p<\infty$, is Grothendieck.
\end{theorem}

\begin{proof}
Apply Lemma \ref{lem:corners} on each of the finitely many central summands
in (\ref{eq:decomp}).  This identifies $M_\psi(\M,\tau)$ with a finite
product of scalar Marcinkiewicz spaces
$M_\psi(\Omega_r,\nu_r)$, one for each triple $(r,i,j)$.

If $\tau(1)<\infty$, every $\nu_r(\Omega_r)$ is finite and (B) gives the
Grothendieck property of every scalar factor by Theorem 6.3 of \cite{DPS}.
If $\tau(1)=\infty$, a component has either finite or infinite measure.  The
zero-end condition in (A) gives (B) on each finite component, while both
parts of (A) apply to every infinite component.  Again every scalar factor
is Grothendieck.  A finite product of Grothendieck spaces is Grothendieck:
its dual is the corresponding finite product of the duals, so weak-star and
weak convergence can be checked coordinatewise.  Banach isomorphisms preserve
the property, proving the first assertion.

Now take the homogeneous atomless model.  Sufficiency was just proved.  The
map $P(x)=e_{11}xe_{11}$ is a bounded projection whose range is isometric to
the scalar corner $M_\psi(0,\gamma)$.  Every complemented subspace of a
Grothendieck space is Grothendieck.  Therefore, if the noncommutative space is
Grothendieck, then so is $M_\psi(0,\gamma)$; Theorem 5.6 of \cite{DPS}
forces (A) or (B).  Finally, for $\psi(t)=t^{1-1/p}$ both endpoint ratios are
$2^{1-1/p}>1$.
\end{proof}

\section{Verification and limitations}
The accompanying script tests (\ref{eq:corners}) on random finite atomic
matrix fields using the discrete Marcinkiewicz norm and verifies exact
reconstruction from entries.  This is a sanity check; the proof is exact.

The finite-degree hypothesis is essential to the present method: the inverse
entry-map bound grows as $r^2$.  For unbounded type-I degree or type-II
algebras, the argument would require a dimension-free decomposition.  The
source proof also cannot simply be copied, because a noncommutative symmetric
space is not a Banach lattice and lacks the suprema and dual lattice
homomorphisms used in Lotz's disjointization argument.  Eight focused upgrade
attempts, recorded in the run attempt note, did not remove these obstacles.

A bounded search through 12 August 2026 checked the run indexes; exact source
wording and title/author searches; and combinations of ``noncommutative
Marcinkiewicz'', ``Grothendieck property'', ``symmetric operator space'',
``matrix amplification'', and ``subhomogeneous''.  It found the source,
adjacent work on noncommutative $L_p+L_q$ spaces and on weak Grothendieck
compactness, but no general resolution or the finite-degree theorem above.
This is not a claim of priority.

\paragraph{Human-review recommendation.}
Check that the scalar corner norm is exactly the zero-extension
Marcinkiewicz norm and verify the standard central decomposition convention
in (\ref{eq:decomp}).  The matrix inequalities and Grothendieck transfer are
then immediate.

\begin{thebibliography}{9}
\bibitem{DPS}
B. de Pagter and F. A. Sukochev,
\emph{The Grothendieck property in Marcinkiewicz spaces},
Indagationes Mathematicae 31 (2020), 791--808;
\href{https://arxiv.org/abs/1912.01162}{arXiv:1912.01162}.
\end{thebibliography}
\end{document}

